\cotefile{10}
\scanpage{77}
% Archive pencil foliation: 76.
quotient de $\handwritten{\overline}{\handwritten{\underline}{Q}}^{\handwritten{*}}$ par le sous-groupe des unités, donc canoniquement isomorphe à $\handwritten{\underline}{Q}$, par l'isomorphisme
\diagram{\begin{tikzcd}[column sep=large]\handwritten{\overline}{\handwritten{\underline}{Q}}^{\handwritten{*}}/\text{unités}\arrow[r]&\handwritten{\underline}{Q}\end{tikzcd}}
donnés par la valuation p-adique $v_p$, normalisée pa r la condition
\[
v_p(p)=1\quad .
\]
\unsure{\deleted{A}L|\deleted{L}A|A|L}'opération de $\deleted{\unsure{\pi|\Pi}_1(\underline{\unsure{\deleted{F}K|\deleted{K}F|K|F|k}}}$ $\operatorname{Gal}(\handwritten{\overline}{\handwritten{\underline}{Q}}_p/\handwritten{\underline}{Q}_p)$ sur ce groupe quotient $Q$ est trivial\unsure{\deleted{i}e|\deleted{e}i|\.{e}|e}, de sorte que $H$ est \handwritten{\uline}{diagonalisable} de groupe des caractères $\handwritten{\underline}{Q}$ :
\[
H=D(\handwritten{\underline}{Q})\quad .
\]
L'augmentation de Tate de $G$ correspond à l'homomorphisme
\diagram{\begin{tikzcd}\handwritten{\underline}{Z}\arrow[r]&\handwritten{\overline}{\handwritten{\underline}{Q}}^{\handwritten{*}},\end{tikzcd}\quad\begin{tikzcd}1\arrow[r,mapsto]&p\quad,\end{tikzcd}}
et l'augmentation induite sur $H$ correspond à l'\handwritten{\uline}{inclusion} canonique
\diagram{\begin{tikzcd}\handwritten{\underline}{Z}\arrow[r,hook]&\handwritten{\underline}{Q}\end{tikzcd}}
(grâee à notre choix de la valuation \handwritten{\uline}{normalisée} comme indiqué\unsure{\deleted{i}e|\deleted{e}i|\.{e}|e}).

\unsure{\rlap{8}B|B|8}.2. Il faut maintenant expliciter l'homomorphisme de \unsure{\deleted{e}g|\deleted{g}e|g}erbes
\diagram{\handwritten{(8\unsure{,|.}2.1)}\qquad\begin{tikzcd}\handwritten{i:}\ \handwritten{\underline}{G}(\handwritten{\underline}{F}_p)\arrow[r]&\handwritten{\underline{\pi}}_p(\handwritten{\underline}{F}_p)\end{tikzcd}}
correspondant à l'inclusion
\diagram{\begin{tikzcd}i:\quad\handwritten{\underline}{Q}_p\text{-faisceaux constructibles}\arrow[r,hook]&\operatorname{Fcriso}(\handwritten{\underline}{F}_p)\end{tikzcd}}
de 5.2. On a vu que l'image essentielle est exactement égale à la sous-catégorie \inserted{pleine} $\unsure{j|J}(\handwritten{\underline}{Q}_p\text{-faisceaux constructibles})$, \deleted{qui est} de sorte que $i$ s'identifie à une\deleted{foncteur} $\handwritten{\unsure{\otimes|\deleted{I}\otimes|\deleted{\otimes}I}}$-équivalence de la catégorie des $\handwritten{\underline}{Q}_p$-faisceaux constructibles avec elle-même. On constate alors que l'homomorphisme correspondant sur les liens $\deleted{\unsure{H\ H|H\mathbin{--}H|H\to H}}$ \begin{tikzcd}[cramped]G'\arrow[r]&G'\end{tikzcd} est l'identité, \deleted{et que} de sorte que $i$ est défini par un torseur sous $G'$, et que ce torseur sous $G'$ peut se décrire ainsi (aurait pu passer à la fin de 7, dans le cas d'un $k$ quelconque ...). On a un homomorphisme canonique
\[
\handwritten{\rlap{\text{(8.2.2)}}\text{(8.2.2)}}\qquad\begin{tikzcd}[cramped]\handwritten{\pi}_1(\handwritten{\underline}{F}_{\substack{\unsure{u.|l.}\\p}}^{\deleted{\handwritten{\unsure{I|1|*|\times}}}})=\overset{\deleted{\handwritten{\unsure{2|\ell|*|\times}}}}{\handwritten{\hat}{\handwritten{\underline}{Z}}}\arrow[r]&G'(\handwritten{\underline}{Q}_p)\end{tikzcd},\quad \handwritten{\text{i.e. }}\begin{tikzcd}[cramped](\pi_1(\underline{F}_p))_{\mathbb{Q}_p}\arrow[r]&G'\end{tikzcd}

\]
d'où un \deleted{homomorphi} foncteur
\diagram{\begin{tikzcd}\text{torseurs}^{\deleted{\underline{Q}_p\unsure{\text{-}|n|r}}}\text{ sous }(\handwritten{\hat}{\handwritten{\underline}{Z}})_{\handwritten{\underline}{Q}_p}\arrow[r]&\text{torseurs sous }G'\ \deleted{\handwritten{\unsure{\varepsilon|e|c|\xi|,|.}}}\end{tikzcd}}

